El Objetivo 3 se desarrolla en tres partes que siguen las fases del enfoque RAD: el desarrollo del sistema y del agente, con el diseño junto al usuario y la construcción; su implementación en el servidor y la transición con el personal; y la evaluación del tiempo de gestión administrativa, que repite la medición inicial con el mismo instrumento (Figura~\ref{fig:ob3_secuencia}). El desarrollo parte de los requerimientos y del diseño ya presentados, que no se repiten.

\begin{figure}[H]
\centering
\begin{tikzpicture}[
  font=\footnotesize,
  fase/.style={draw, rounded corners=2pt, line width=0.6pt, align=center,
               minimum height=1.7cm, text width=2.4cm, inner sep=3pt},
  flecha/.style={-{Stealth[length=2.2mm]}, line width=0.7pt, draw=dExternoB},
  obj/.style={font=\scriptsize\bfseries, text=dExternoB}
]
\node[fase, fill=dPersonaF, draw=dPersonaB] (diag) {\textbf{Diagnóstico}\\ $TGA$ pretest\\ RF y RNF};
\node[fase, fill=dSistemaF, draw=dSistemaB, right=0.25cm of diag] (dis) {\textbf{Diseño}\\ arquitectura, datos e interfaz};
\node[fase, fill=dAgenteF, draw=dAgenteB, right=0.25cm of dis] (des) {\textbf{Desarrollo}\\ prototipo y tarjetas Kanban};
\node[fase, fill=dAgenteF, draw=dAgenteB, right=0.25cm of des] (imp) {\textbf{Implementación}\\ servidor, WhatsApp y puesta en marcha};
\node[fase, fill=dClinicoF, draw=dClinicoB, right=0.25cm of imp] (eval) {\textbf{Evaluación}\\ $TGA$ postest};
\foreach \a/\b in {diag/dis, dis/des, des/imp, imp/eval} \draw[flecha] (\a) -- (\b);
\node[obj, above=2pt of diag] {Objetivo 1};
\node[obj, above=2pt of dis] {Objetivo 2};
\node[obj, above=2pt of imp] {Objetivo 3};
\draw[flecha, dashed] (eval.south) -- ++(0,-0.45) -| node[pos=0.25, below, font=\scriptsize]{comparación pretest--postest} (diag.south);
\end{tikzpicture}
\caption{Secuencia de los objetivos de la investigación}
\label{fig:ob3_secuencia}
\end{figure}
